\chapter{Neural ODEs, CDEs, and SDEs}
\label{ch:nodes}

\epigraph{\itshape A residual block is a discrete Euler step. Take the
limit and you have an ODE.}{Chen, Rubanova, Bettencourt, Duvenaud (2018)}

\section{From ResNets to Neural ODEs}
A residual block computes
$\bm{h}_{n+1}=\bm{h}_n+\bm{f}_\theta(\bm{h}_n)$. Treating depth as time
and taking $\Delta t\to 0$ gives the \emph{neural ODE}~\cite{chen2018neuralode}:
\begin{equation}
\frac{\dd \bm{h}(t)}{\dd t}=\bm{f}_\theta(\bm{h}(t),t),\qquad
\bm{h}(0)=\bm{x}.
\end{equation}
The ``forward pass'' is an ODE solve from $0$ to $T$. Memory cost is
$\Oh(1)$ in the number of steps because gradients are computed by
solving a backward adjoint ODE rather than storing intermediate states.

\subsection{Adjoint method}
Define the loss $L(\bm{h}(T))$. Introduce the adjoint
$\bm{a}(t)=\partial L/\partial \bm{h}(t)$, which satisfies
\begin{equation}
\frac{\dd \bm{a}(t)}{\dd t}=-\bm{a}(t)^\T \frac{\partial \bm{f}_\theta}{\partial \bm{h}}.
\end{equation}
The parameter gradient is
$\dd L/\dd\theta=-\int_0^T \bm{a}(t)^\T (\partial \bm{f}_\theta/\partial\theta)\,\dd t$.
This is the continuous-time analogue of backpropagation.

\begin{pitfall}{Adjoint instability}
The continuous adjoint can be numerically inaccurate when the forward
solver uses adaptive step control or interpolated dense output. The
\emph{interpolated} or \emph{checkpointed} adjoint, and JAX/PyTorch
\texttt{torchdiffeq} backsolve, mitigate this. For chaotic systems,
adjoints diverge along Lyapunov directions; consider shadowing /
ensemble adjoints.
\end{pitfall}

\section{Variants and architectures}
\paragraph{Augmented neural ODEs}
The flow of an ODE is a homeomorphism, so a neural ODE cannot represent
maps that change topology. Augmenting the state with extra dimensions
$\bm{a}(0)=\bm{0}$ restores expressivity~\cite{dupont2019anode}.

\paragraph{Latent ODEs and continuous normalising flows.}
For irregular time series and density estimation, model dynamics in a
latent space and decode~\cite{rubanova2019latentode}. Continuous
normalising flows (CNFs) use the instantaneous change-of-variables
formula
$\dd\log p/\dd t=-\Tr(\partial \bm{f}/\partial \bm{h})$ to evaluate
likelihoods.

\paragraph{Neural CDEs.}
For irregular time series, treat the data path $\bm{X}(t)$ as a
control and integrate
$\dd\bm{h}(t)=\bm{f}_\theta(\bm{h}(t))\,\dd\bm{X}(t)$ in the sense of
controlled differential equations~\cite{kidger2020ncde}.

\paragraph{Neural SDEs.}
Add Brownian motion: $\dd\bm{h}=\bm{f}_\theta\,\dd t+\bm{\sigma}_\theta\,\dd\bm{W}_t$.
Trainable via Wasserstein matching, KL with reference SDE, or score-matching;
foundational for generative diffusion (\cref{ch:generative}).

\section{Hamiltonian and Lagrangian neural networks}
Architectural priors that respect mechanics:
\begin{itemize}
\item \textbf{HNN}~\cite{greydanus2019hnn}: parameterise $H_\theta(\bm{q},\bm{p})$;
integrate $\dot{\bm{q}}=\partial_{\bm{p}}H_\theta$,
$\dot{\bm{p}}=-\partial_{\bm{q}}H_\theta$.
\item \textbf{LNN}~\cite{cranmer2020lnn}: parameterise $L_\theta(\bm{q},\dot{\bm{q}})$
and impose Euler--Lagrange equations.
\item \textbf{Symplectic ODE-Net, Constrained HNN}: enforce constraints
and conservation explicitly.
\end{itemize}
These networks generalise far beyond their training horizon precisely
because they encode the right invariant.

\section{ODE-based deep generative models}
\subsection{Continuous normalising flows (CNF)}
Train $\bm{f}_\theta$ so that the pushforward of a base density
matches the data density:
\begin{equation}
\log p_\theta(\bm{x})=\log p_0(\bm{h}(0)) - \int_0^T \nabla\!\cdot\bm{f}_\theta(\bm{h}(t),t)\,\dd t.
\end{equation}
\subsection{FFJORD and Hutchinson trace estimators}
Replace the exact trace by a Hutchinson estimator
$\Tr(\bm{J})\approx \E_{\bm{\epsilon}\sim\mathcal{N}}[\bm{\epsilon}^\T\bm{J}\bm{\epsilon}]$ for tractable training.

\subsection{Flow matching and stochastic interpolants}
Recent advances~\cite{lipman2023flowmatching,albergo2023interpolants}
train CNFs by regressing onto a target velocity field defined by a
prescribed interpolation between base and data distributions, avoiding
expensive ODE simulation during training. Discussed at length in
\cref{ch:generative}.

\section{Practical training tips}
\begin{recipe}{Training a neural ODE end-to-end}
\begin{enumerate}
\item Normalise inputs/outputs; centre time near zero.
\item Start with a permissive solver (Tsit5, atol/rtol $10^{-3}$) and
tighten as training progresses.
\item Use the \emph{checkpointed adjoint} for long horizons; the
\emph{recursive checkpoint} for memory.
\item Add Jacobian/Frobenius penalties to discourage stiffness.
\item Curriculum on $T$: train on short windows, extend gradually.
\end{enumerate}
\end{recipe}

\section{Where neural ODEs shine and where they struggle}
\textbf{Strengths:} natural for irregular time series, continuous-depth
generative modelling, mechanistic priors via HNN/LNN, smooth and
invertible models.\\
\textbf{Weaknesses:} high wall-clock cost from solver calls;
representation limited to homeomorphisms (without augmentation);
delicate adjoint behaviour on stiff or chaotic dynamics.

\begin{insight}{Neural ODE = parametric differential equation}
The most productive way to think about neural ODEs is not as a deep
network whose depth is continuous, but as a \emph{parametric ODE family
fitted to data}. This perspective unifies them with Universal
Differential Equations (\cref{ch:hybrid-ude}), system identification
(\cref{ch:operator-theoretic}), and operator learning
(\cref{ch:operators}).
\end{insight}
